\documentclass[11pt,a4paper]{article}
\newcommand{\intestazione}{Trimero ad anello --- Schema dei passi sotto rumore}
% =====================================================================
%  Preambolo comune ai documenti sul dimero (serie "che cosa abbiamo fatto")
%  Incluso con % =====================================================================
%  Preambolo comune ai documenti sul dimero (serie "che cosa abbiamo fatto")
%  Incluso con % =====================================================================
%  Preambolo comune ai documenti sul dimero (serie "che cosa abbiamo fatto")
%  Incluso con \input{preambolo_dimero} subito dopo \documentclass.
% =====================================================================
\usepackage[T1]{fontenc}
\usepackage[utf8]{inputenc}
\usepackage[main=italian,provide=*]{babel}
\usepackage{amsmath,amssymb,mathtools}
\usepackage{graphicx}
\usepackage{booktabs}
\usepackage{colortbl}
\usepackage{xcolor}
\usepackage{geometry}
\usepackage{placeins}
\usepackage{caption}
\usepackage{enumitem}
\usepackage{fancyhdr}
\usepackage{needspace}
\usepackage{tikz}
\usetikzlibrary{arrows.meta,positioning,shapes.geometric,calc,fit,backgrounds}
\usepackage{hyperref}
\hypersetup{colorlinks=true,linkcolor=blu,citecolor=blu,urlcolor=blu}

\geometry{margin=2.5cm}

% --- colori coerenti con quelli delle figure (palette Okabe-Ito)
\definecolor{blu}{HTML}{0072B2}
\definecolor{arancio}{HTML}{E69F00}
\definecolor{verdes}{HTML}{009E73}
\definecolor{rossos}{HTML}{D55E00}
\definecolor{violas}{HTML}{CC79A7}
\definecolor{grigetto}{HTML}{E8E8E8}

\captionsetup{font=small,labelfont=bf,width=0.92\textwidth}

\pagestyle{fancy}
\fancyhf{}
\fancyhead[L]{\small\itshape\intestazione}
\fancyhead[R]{\small\thepage}
\renewcommand{\headrulewidth}{0.3pt}

% --- notazione
\newcommand{\ket}[1]{\lvert #1\rangle}
\newcommand{\bra}[1]{\langle #1 \rvert}
\newcommand{\media}[1]{\langle #1 \rangle}
\newcommand{\vs}[1]{\mathbf{s}_{#1}}

% --- riquadro "in una riga": il messaggio da portare a casa
\newcommand{\messaggio}[1]{%
  \begin{center}
  \begin{minipage}{0.93\textwidth}
  \setlength{\fboxsep}{9pt}%
  \colorbox{grigetto}{\begin{minipage}{\dimexpr\textwidth-2\fboxsep}
  \small\textbf{In una riga.}\enspace #1
  \end{minipage}}
  \end{minipage}
  \end{center}
}

% --- riquadro "come lo sappiamo": la verifica numerica dietro un'affermazione
\newcommand{\verifica}[1]{%
  \begin{center}
  \begin{minipage}{0.93\textwidth}
  \small\textcolor{verdes}{\rule{2.2em}{1.1pt}}\enspace
  \textcolor{verdes}{\textbf{Come lo sappiamo.}}\enspace #1
  \end{minipage}
  \end{center}
  \medskip
}
 subito dopo \documentclass.
% =====================================================================
\usepackage[T1]{fontenc}
\usepackage[utf8]{inputenc}
\usepackage[main=italian,provide=*]{babel}
\usepackage{amsmath,amssymb,mathtools}
\usepackage{graphicx}
\usepackage{booktabs}
\usepackage{colortbl}
\usepackage{xcolor}
\usepackage{geometry}
\usepackage{placeins}
\usepackage{caption}
\usepackage{enumitem}
\usepackage{fancyhdr}
\usepackage{needspace}
\usepackage{tikz}
\usetikzlibrary{arrows.meta,positioning,shapes.geometric,calc,fit,backgrounds}
\usepackage{hyperref}
\hypersetup{colorlinks=true,linkcolor=blu,citecolor=blu,urlcolor=blu}

\geometry{margin=2.5cm}

% --- colori coerenti con quelli delle figure (palette Okabe-Ito)
\definecolor{blu}{HTML}{0072B2}
\definecolor{arancio}{HTML}{E69F00}
\definecolor{verdes}{HTML}{009E73}
\definecolor{rossos}{HTML}{D55E00}
\definecolor{violas}{HTML}{CC79A7}
\definecolor{grigetto}{HTML}{E8E8E8}

\captionsetup{font=small,labelfont=bf,width=0.92\textwidth}

\pagestyle{fancy}
\fancyhf{}
\fancyhead[L]{\small\itshape\intestazione}
\fancyhead[R]{\small\thepage}
\renewcommand{\headrulewidth}{0.3pt}

% --- notazione
\newcommand{\ket}[1]{\lvert #1\rangle}
\newcommand{\bra}[1]{\langle #1 \rvert}
\newcommand{\media}[1]{\langle #1 \rangle}
\newcommand{\vs}[1]{\mathbf{s}_{#1}}

% --- riquadro "in una riga": il messaggio da portare a casa
\newcommand{\messaggio}[1]{%
  \begin{center}
  \begin{minipage}{0.93\textwidth}
  \setlength{\fboxsep}{9pt}%
  \colorbox{grigetto}{\begin{minipage}{\dimexpr\textwidth-2\fboxsep}
  \small\textbf{In una riga.}\enspace #1
  \end{minipage}}
  \end{minipage}
  \end{center}
}

% --- riquadro "come lo sappiamo": la verifica numerica dietro un'affermazione
\newcommand{\verifica}[1]{%
  \begin{center}
  \begin{minipage}{0.93\textwidth}
  \small\textcolor{verdes}{\rule{2.2em}{1.1pt}}\enspace
  \textcolor{verdes}{\textbf{Come lo sappiamo.}}\enspace #1
  \end{minipage}
  \end{center}
  \medskip
}
 subito dopo \documentclass.
% =====================================================================
\usepackage[T1]{fontenc}
\usepackage[utf8]{inputenc}
\usepackage[main=italian,provide=*]{babel}
\usepackage{amsmath,amssymb,mathtools}
\usepackage{graphicx}
\usepackage{booktabs}
\usepackage{colortbl}
\usepackage{xcolor}
\usepackage{geometry}
\usepackage{placeins}
\usepackage{caption}
\usepackage{enumitem}
\usepackage{fancyhdr}
\usepackage{needspace}
\usepackage{tikz}
\usetikzlibrary{arrows.meta,positioning,shapes.geometric,calc,fit,backgrounds}
\usepackage{hyperref}
\hypersetup{colorlinks=true,linkcolor=blu,citecolor=blu,urlcolor=blu}

\geometry{margin=2.5cm}

% --- colori coerenti con quelli delle figure (palette Okabe-Ito)
\definecolor{blu}{HTML}{0072B2}
\definecolor{arancio}{HTML}{E69F00}
\definecolor{verdes}{HTML}{009E73}
\definecolor{rossos}{HTML}{D55E00}
\definecolor{violas}{HTML}{CC79A7}
\definecolor{grigetto}{HTML}{E8E8E8}

\captionsetup{font=small,labelfont=bf,width=0.92\textwidth}

\pagestyle{fancy}
\fancyhf{}
\fancyhead[L]{\small\itshape\intestazione}
\fancyhead[R]{\small\thepage}
\renewcommand{\headrulewidth}{0.3pt}

% --- notazione
\newcommand{\ket}[1]{\lvert #1\rangle}
\newcommand{\bra}[1]{\langle #1 \rvert}
\newcommand{\media}[1]{\langle #1 \rangle}
\newcommand{\vs}[1]{\mathbf{s}_{#1}}

% --- riquadro "in una riga": il messaggio da portare a casa
\newcommand{\messaggio}[1]{%
  \begin{center}
  \begin{minipage}{0.93\textwidth}
  \setlength{\fboxsep}{9pt}%
  \colorbox{grigetto}{\begin{minipage}{\dimexpr\textwidth-2\fboxsep}
  \small\textbf{In una riga.}\enspace #1
  \end{minipage}}
  \end{minipage}
  \end{center}
}

% --- riquadro "come lo sappiamo": la verifica numerica dietro un'affermazione
\newcommand{\verifica}[1]{%
  \begin{center}
  \begin{minipage}{0.93\textwidth}
  \small\textcolor{verdes}{\rule{2.2em}{1.1pt}}\enspace
  \textcolor{verdes}{\textbf{Come lo sappiamo.}}\enspace #1
  \end{minipage}
  \end{center}
  \medskip
}

\usepackage{pdflscape}
\usepackage{array}

\title{\bfseries Schema riassuntivo --- Parte 2\\[2pt]
\large Punti di lavoro, risultati chiave e collegamenti sotto rumore, trimero ad anello}
\author{Samuele Schiavi \\ \small Tesi triennale in Fisica, Università di Parma
--- Relatore: Prof.\ Alessandro Chiesa}
\date{}

\begin{document}
\maketitle
\thispagestyle{fancy}

Questo documento non introduce fisica nuova: raccoglie in un unico posto,
per ciascuno dei sette documenti già scritti (cinque stadi più due
estensioni facoltative), il punto di lavoro usato, il risultato chiave e
il collegamento con quanto viene dopo. Mirror diretto dello schema
analogo della Parte 2 sul dimero.

\section*{Un confronto col dimero da tenere presente}

Il chiarimento aperto sul dimero (``$N^*$ massimizza il segnale
misurabile, non l'accuratezza della stima --- $|C(N^*)|$ può sovrastimare
il vero $C(t)$ anche di un fattore $2\times$'') \textbf{non si ripete
identico qui}: per il correlatore e il punto di lavoro di questa serie, a
rumore nullo $N^*$ dà un valore praticamente coincidente col vero
$|C_\text{esatto}(t{=}2)|$ (Documento 8). Non è una regola generale che si
trasporta da un sistema all'altro --- va verificata di nuovo ogni volta,
come qui.

\section{Tabella d'insieme --- i cinque stadi}

\renewcommand{\arraystretch}{1.3}
\begin{landscape}
\thispagestyle{empty}
\begin{center}
\small
\setlength{\tabcolsep}{3.5pt}
\begin{tabular}{@{}>{\bfseries\raggedright}p{2.3cm} p{4.1cm} p{4.1cm} p{4.1cm} p{4.1cm}@{}}
\toprule
\rowcolor{grigetto}
 & \textbf{5. Il modello} & \textbf{6. Il VQE} & \textbf{7. La dinamica} & \textbf{8. I correlatori} \\
\midrule
Obiettivo &
Due canali (depolarizzante 1q/2q, readout), calibrati su hardware reale,
riusati dal dimero &
Degradazione del VQE ideale (Parte 1) eseguito su simulatore rumoroso &
$N^*$: compromesso Trotter/rumore, due scenari (A e $R_0$) &
$N^*$ del correlatore, readout genuinamente campionato \\
\addlinespace
Punto/i di lavoro &
$\varepsilon_{1q}{=}2.9{\times}10^{-4}$, $\varepsilon_{2q}{=}3.8{\times}10^{-3}$,
$p_\text{readout}{=}2.3{\times}10^{-2}$ (\texttt{ibm\_torino}) &
Scenario A ($b_c{=}2.4$, $D{=}0.15$) &
Scenario A \emph{e} $R_0$ ($b{=}0.05$, $D{=}1.93$) &
Scenario A, $C_{21}^{xx}$ \\
\addlinespace
Risultato chiave &
identico al dimero, verificato indipendente dal numero di qubit &
$1{-}F{=}2.98\%$ ($2\times$ il dimero); costo preparazione $6$ CNOT (il
più caro della pipeline, non il più economico) &
$N^*_A=3$; $N^*_B=9$ (correzione di una griglia troppo rada, poi
propagata) &
$N^*=3$; a rumore nullo $|C(N^*)|$ quasi coincide col vero valore
(diverso dal dimero) \\
\addlinespace
Collegamento &
fornisce il \texttt{NoiseModel} riusato in tutta la Parte 2 &
fissa la scala di degradazione di base &
introduce $N^*$ per entrambi gli scenari, ripreso nel Documento 9 &
fornisce baseline e correlatore per i Documenti 9--11 \\
\bottomrule
\end{tabular}

\vspace{6pt}
\begin{tabular}{@{}>{\bfseries\raggedright}p{2.3cm} p{16.9cm}@{}}
\toprule
\rowcolor{grigetto}
 & \textbf{9. Lo scan sui parametri} \\
\midrule
Obiettivo & Come si sposta $N^*$ al variare di $\varepsilon_{1q},\varepsilon_{2q},p_\text{readout}$ \\
\addlinespace
Punto/i di lavoro & Scenario A ($b/J{=}2.4$, $D/J{=}0.15$); scan indipendente su ciascun parametro, altri fissati al riferimento \\
\addlinespace
Risultato chiave & $N^*$ Trotter monotono non-crescente in $\varepsilon_{2q}$ ($4\to2$); $N^*$ correlatore fermo a $3$ su tutto il range --- differenza fra le due metriche non vista in questa forma sul dimero; invarianza da $p_\text{readout}$ confermata a shot finiti \\
\addlinespace
Collegamento & chiude i cinque stadi; apre le due estensioni facoltative (Documenti 10, 11) \\
\bottomrule
\end{tabular}
\end{center}
\end{landscape}

\clearpage

\section{Tabella d'insieme --- le due estensioni}

\begin{center}
\small
\renewcommand{\arraystretch}{1.35}
\begin{tabular}{@{}p{3.4cm}p{6.6cm}p{6.6cm}@{}}
\toprule
\rowcolor{grigetto}
 & \textbf{10. VQE noise-aware} & \textbf{11. Readout asimmetrico} \\
\midrule
Domanda &
Vale la pena riottimizzare vedendo il rumore? &
Cosa succede se $p_{01}\neq p_{10}$? \\
\addlinespace
Esito &
\textbf{Diverso dal dimero}: energia migliore, fedeltà peggiore
($\Delta E{=}{-}4.16{\times}10^{-3}$, $\Delta F{=}{-}1.57{\times}10^{-2}$,
a struttura di circuito fissa); $N^*$ comunque invariato su entrambe le
metriche &
\textbf{Come sul dimero}: nessuna sorpresa, $N^*=3$ invariato dal
simmetrico fino a rapporti $50\times$ \\
\addlinespace
Punto notevole &
Isolato un artefatto del transpilatore (riduzione CNOT a valori speciali
di $\theta$) mai visto sul dimero, distinto dal risultato fisico &
Verificato che il registro a 3 qubit (contro i 2 del dimero) non
complica la correzione: si applica a un solo bit misurato \\
\bottomrule
\end{tabular}
\end{center}

\section{Il modello di rumore}

Due canali, calibrazione \texttt{ibm\_torino}, letteralmente lo stesso
modello del dimero (wrapper, non una riscrittura) --- verificato
indipendente dal numero di qubit per costruzione dell'API di Qiskit.

\section{Il VQE sotto rumore}

Riuso dei parametri ideali della Parte 1, nessuna riottimizzazione:
$1-F=2.98\%$, il doppio del dimero a parità di rumore di riferimento,
coerente col numero di CNOT triplo ($6$ contro $2$).

\messaggio{L'ansatz dell'anello è il più costoso in gate di tutta la
pipeline, non il più economico come sul dimero --- una differenza
strutturale che si ripercuote su ogni stadio successivo.}

\section{La dinamica sotto rumore}

Due scenari trattati esplicitamente (a differenza del dimero, un solo
punto): $N^*_A=3$ (preparazione VQE), $N^*_B=9$ (da $\ket{000}$, punto
$R_0$). Nessun readout coinvolto --- fedeltà calcolata direttamente
sull'operatore densità, numeri identici a quelli già stabiliti prima di
questa serie narrativa.

\section{I correlatori sotto rumore}

Passaggio al readout genuinamente campionato (shot finiti), come già
fatto sul dimero in una revisione successiva alla prima stesura di questa
pipeline. $N^*=3$ confermato. La scoperta di questo documento: a
differenza del dimero, qui $N^*$ non sovrastima il vero correlatore
fisico --- un fenomeno specifico del sistema e del punto, non universale.

\section{Lo scan sui parametri di rumore}

Monotonia confermata per la fedeltà di Trotter; il correlatore mostra
invece un $N^*$ sorprendentemente stabile su tutto il range di
$\varepsilon_{2q}$ esplorato --- osservazione aperta, non ancora spiegata.

\section{Le due estensioni}

L'estensione VQE noise-aware (Documento 10) è quella con l'esito più
diverso dal dimero di tutta questa Parte 2: un compromesso energia/fedeltà
reale, isolato da un artefatto genuino del transpilatore specifico
dell'ansatz a sei parametri. L'estensione readout asimmetrico (Documento
11) chiude invece senza sorprese, esattamente come sul dimero.

\messaggio{Due estensioni, due esiti opposti: una conferma che ``fare la
stessa cosa'' su un sistema diverso non garantisce lo stesso risultato,
anche quando l'argomento teorico di partenza è identico.}

\end{document}
